\cotefile{156-7}
\scanpage{26}\mspage{25}

D'abord, \deleted{\unsure{prenons|posons}} dégageons quelques conditions concernant $\leq$, $\ll$, \deleted{\unsure{et}} indépendamment \unsure{de|des} $\parallel$.

\[
X\leq Y\ \Longrightarrow\ X\ll Y \qquad \text{pour mémoire}
\tag*{(ML \deleted{\unsure{1}}0)}
\]

(ML 1) Soit
\diagram{\begin{tikzcd}[ampersand replacement=\&,row sep=0pt,column sep=3pt]
X' \& {\leq} \& X\\
\& \& {\unsure{\rotatebox{90}{$\leq$}|\rotatebox{90}{$\ll$}|\Downarrow}}\\
\& \& X''
\end{tikzcd}}
, soit $x\in L$\inserted{[on doit \unsure{pouvoir} remplacer $x$ par un $Z\in M$ quelc.?]} avec $x\ll X'$, $x\ll X''$, alors $\exists Z$ \deleted{\unsure{$\in$}} avec $x\ll Z$ et $Z\leq X',X''$.
% slanted
\margin{$\Big\Vert$ Il faudra \unsure{renforcer}, avec $X'\ll X$ $X''\leq X$, par les propriétés \unsure{du|de} $\mathrm{Inf}^{\ll}(X',X'')$ $\Big\}$}

[I.e \deleted{$X'\cap X''$} \unsure{la} ``figure\unsure{''} ``intersection'' $X'\cap X''=\widetilde X'\cap\widetilde X''$ est un $\mathrm{Inf}$ de $X', X''$ pour $\ll$, du moins \unsure{p.\,v.\,à|p.\,r.\,à|p. v. :} $L$)

\deleted{\unsure{(ML 2)} \unsure{$\forall$}} Posons pour \deleted{\unsure{$\mathrm{omb}(X)\mathrel{=}\mathrm{omb}$}}\inserted{$X\in M$}
\[
\begin{aligned}
\mathrm{omb}(X)&=\{x\in L\mid x\ll X\}\\
\mathrm{omb}(X)^\circ&=\mathrm{omb}(X)\smallsetminus\bigcup_{Y<X}\mathrm{omb}(Y)
\end{aligned}
\]

Posons
\[
\begin{aligned}
x\overset{\circ}{\ll}X &\ \overset{\mathrm{def}}{\Longleftrightarrow}\ x\in\mathrm{omb}(X)^\circ \qquad \text{i.e.}\\
&\qquad x\ll X,\ \text{et}\ x\ll X'\leq X\ \Longrightarrow\ X'=X
\end{aligned}
\]

On définit de $=$
\[
\begin{aligned}
Y\overset{\circ}{\ll}X &\ \Longleftrightarrow\ Y\ll X,\ \text{et}\ Y\ll X'\leq X\ \Longrightarrow\ X'=X\\
&\qquad \text{i.e. }X\text{ minimal dans }\widetilde X^{Y}=\{X'\in\widetilde X\mid \deleted{\unsure{Y\ll}}Y\ll X'\}
\end{aligned}
\]

\deleted{(ML 2) \unsure{$\forall$} Soit $X\in M$. Alors $\mathrm{omb}(X)^\circ\neq\varnothing$}

\deleted{(ML 3) Soit $X\in M$, $x\in L$, $x\ll\unsure{X|1|l}$}
\deleted{\unsure{$\exists Z,\ X^{\circ},\,X'\leq X.$}}

% slanted
\margin{\deleted{\unsure{ML 1}} (ML 2) \uline{avant} le lemme 1 $\Big\}$}

\uline{Lemme 1} Soient $\deleted{\unsure{X'}}X,\deleted{\unsure{Y'}}Y\in M$, \deleted{\unsure{avec}} \unsure{Supposons|Suppos} \unsure{conditions équiv.}
\[
\begin{aligned}
&\text{(i)} && X\leq Y\\
&\text{(ii)} && X\ll Y\\
&\text{(ii\unsure{bis})} && \mathrm{omb}\,X\subset\mathrm{omb}\,Y
     && \deleted{\Downarrow\,\text{\unsure{via|par} (ML \unsure{2|3})}}\\
&\text{(iii)} && \mathrm{omb}(X)^\circ\subset\mathrm{omb}(Y)
     && \Downarrow\,\text{via (ML 2)}\\
&\text{(iv)} && \mathrm{omb}(X)^\circ\cap\mathrm{omb}(Y)\neq\varnothing
\end{aligned}
\]

Il suffit de prouver $(\mathrm{iv})\Longrightarrow(\mathrm{i})$ Or si $x\in\mathrm{omb}(X)^\circ\cap\mathrm{omb}(Y)$, \deleted{\unsure{on a $x$} $\exists$} par ML 1 $\exists Z$, $Z\leq\deleted{\unsure{X'}}X,\deleted{\unsure{Y'}}Y$, avec $x\ll Z$\deleted{\unsure{$\ll$}}. On a $x\ll Z\leq X$, d'où $Z=X$ puisque $x\overset{\circ}{\ll}X$, d'où $X\leq Y$ \unsure{qed.}
\scanpage{27}\mspage{26}

\uline{Cor} Si $\mathrm{omb}(X)^\circ\cap\mathrm{omb}(Y)^\circ\neq\varnothing$, on a $X=Y$ / i.e. $X\neq Y\Rightarrow\mathrm{omb}(X)^\circ\cap\mathrm{omb}(Y)^\circ=\varnothing$

\margin{\unsure{$\mathrm{At}\,L\,1|\mathrm{At}_{L}\,1$}}
\[
\left|\begin{aligned}
&\uline{\text{ML 2}}\qquad \forall X\in M,\quad \mathrm{omb}(X)^\circ\neq\varnothing.\quad \text{i.e. }\exists x\in L,\quad x\overset{\circ}{\ll}X\\
&\uline{\text{ML 3}}\qquad \forall X\in M,\quad \text{les }\mathrm{omb}(Y)^\circ\quad(Y\in\widetilde X)\quad\text{recouvrent}\\
&\mathrm{omb}(X),\quad\text{i.e. }\deleted{\unsure{a}}\forall x\in L,\quad x\ll X,\quad\exists\deleted{\unsure{X}}Y\leq X,\ \text{avec }x\overset{\circ}{\ll}Y.
\end{aligned}\right.
\]
% slanted
\margin{$\mathrm{At}_{M}\,4$ a) pour $L$}

% slanted
\margin{ML 2 et ML 3 \deleted{\unsure{les|dès les}} signifient que les $\mathrm{omb}(Y)^\circ$ forment une \uline{partition} de $X$.}

\uline{NB} Par le cor. ci-dessus, cet $Y$ est unique.

Par le lemme, si \deleted{\unsure{$x\in|x\ll$}} $Y'\leq X$, et $x\ll Y'$, comme $x\in\mathrm{omb}(Y)^\circ\cap\mathrm{omb}(Y')$ \deleted{\unsure{alors}} on a $Y\leq Y'$ par le lemme, donc $Y'$ est le plus petit \unsure{élémt|élément} de $\widetilde X^x$. Plus généralement :

\deleted{\uline{Lemme 2} Soit $X,Y\in M$. $Y\ll X$. Alors}
\figure{\begin{tikzpicture}[x=1cm,y=1cm]
\draw (0,.24)--(.20,.24)--(.22,.01);
\draw (0,0)--(.19,0);
\draw (.015,.005)--(.26,.56);
\draw (.085,.005)--(.32,.55);
\end{tikzpicture}}

On n'a pas jusqu'à présent, \unsure{utilisé|utiliser} la définition du \unsure{sous-ens|sous-ens.} $L$ de $M$, comme formé d'él. minimaux pour $\ll$. Donc \unsure{Tout|tout} marche avec un $L$ quelconque, pourvu qu'on suppose ML 2, ML 3 pour cet $L$. Si on \unsure{prenait|prend} $L=M$, ML 2 est trivial, ML 3 signifie que $\forall\deleted{\unsure{x\in}}X\in M$ et $Z\ll X$, $\exists Y\in\widetilde X$ tel que $Z\overset{\circ}{\ll}Y$, i.e. $\mathrm{At}_{M}\,4$ a)\unsure{''} (page 6). Mais il faudra dans ML 1 \deleted{\unsure{supposer que}} \inserted{généraliser}, en remplaçant $x$ par un $Z\inserted{\in M}$ quelconque.

Mais il devrait être possible de le prouver, via ML 3 \inserted{et ML 2} pour $L$, . L'idée est pour $\deleted{\unsure{Y}}Z\ll X$, de prendre $x\in L$, $x\overset{\circ}{\ll}Z$, \inserted{(ML \unsure{2|3})} et l'unique $Y\in\widetilde X$ tel que $x\overset{\circ}{\ll}Y$ (ML 3), et il \deleted{\unsure{po|pe}} faudrait prouver que $\deleted{\unsure{Y}}Z\overset{\circ}{\ll}\deleted{\unsure{Z}}Y$. Ceci se \unsure{décompose|discompose} en deux : si $x\ll Y\leq X$, \unsure{on a} $Z\ll Y$ et alors $x\overset{\circ}{\ll}Y$ \unsure{sss|ssi} $Z\overset{\circ}{\ll}Y$.
\scanpage{28}\mspage{27}

Donc on voudrait prouver

\uline{Lemme} 2 Soit
\diagram{\begin{tikzcd}[ampersand replacement=\&,row sep=0pt,column sep=3pt]
Z \& {\ll} \& X\\
\& \& \rotatebox{90}{$\leq$}\\
\& \& {\deleted{\unsure{Y}}X'}
\end{tikzcd}}
, soit $x\in L$, tel que $x\ll Z$, \deleted{\unsure{$x$}} $x\ll X'$. Si $x\overset{\circ}{\ll}Z$, on aura $Z\ll X'$, et alors $x\overset{\circ}{\ll}X'\Longleftrightarrow Z\overset{\circ}{\ll}X'$.

\uline{On admettra} ML 1 \uline{sous forme généralisée}, pour conclure ici que $\exists Z'\leq Z$, $\deleted{\unsure{\text{avec}|Z'}}Z'\ll X'$ $Z'\ll X'$ et $x\ll Z'$, :
\diagram{\begin{tikzcd}[ampersand replacement=\&,row sep=1pt,column sep=1pt]
\& Z \& {\ll} \& X\\
\& \rotatebox{90}{$\leq$} \& \& \rotatebox{90}{$\leq$}\\
\& Z' \& {\ll} \& X'\\
{\rotatebox{48}{$\ll$}} \& \& \&\\
x \& \& \&
\end{tikzcd}}

Ceci dit, si $x\overset{\circ}{\ll}Z$, on a $Z'=Z$, donc $Z\ll X'$. Si de plus $x\overset{\circ}{\ll}X'$, alors \uline{a fortiori} \deleted{$Z'\overset{\circ}{\ll}\ \unsure{x}\ll$} $Z\overset{\circ}{\ll}X'$ (tautologique), prouvons \unsure{l'inverse|l'envers|l'assertion} \deleted{\unsure{$Z$}}
\[
x\overset{\circ}{\ll}Z\overset{\circ}{\ll}X'\ \Longrightarrow\ x\overset{\circ}{\ll}X'.
\]

Soit en effet \deleted{$X''\leq\unsure{X'|X}$} $X''$ avec $x\ll X''\leq X'$, par ce qui précède appliqué à $X''\leq X$ (\unsure{ald|et à} $X'\leq X$) on a $\deleted{\unsure{x}}Z\deleted{\unsure{\leq}}\ll X''\leq X'$, et comme $Z\overset{\circ}{\ll}X'$ on conclut $X''=X'$, qed.

Ceci prouve donc que $\forall Z$ \deleted{\unsure{$\in$}}tel que $Z\ll X$, $\exists X'\leq X$ avec $Z\overset{\circ}{\ll}X'$ : (en utilisant ML 1 renforcé, ML 2 et ML 3) C'est $\mathrm{At}_{M}\,4$ a). On trouve en $=$ temps
\scanpage{29}\mspage{28}

\deleted{\uline{Corollaire} Soit}
\[
\begin{aligned}
&\deleted{Z\ll Y\ll X}\\
&\deleted{X\ll\unsure{Z|Y}\ll Z\unsure{.}}\qquad .
\end{aligned}
\]
\deleted{Conditions équivalentes}
\[
\deleted{\text{(i) }Z\overset{\circ}{\ll}\unsure{,|Y}\quad\text{et }Z\ll X}
\]

\uuline{Corollaire 1} Pour $Z\ll X$, $\exists X'$ avec $Z\overset{\circ}{\ll}X'\deleted{\unsure{\ll|<}}\leq X$.

\uline{Corollaire 2} Soit
\diagram{\begin{tikzcd}[ampersand replacement=\&,row sep=0pt,column sep=3pt]
Z \& {\ll} \& X\\
\& \& \rotatebox{90}{$\leq$}\\
\& \& X'
\end{tikzcd}}
. Conditions \unsure{équiv.}
\[
\begin{aligned}
&\deleted{\text{(i) }Z\overset{\circ}{\unsure{\ll|\leq}}X}\\
&\deleted{\text{(ii) }\unsure{\mathrm{om}|\mathrm{omb}}(Z)^\circ\subset\mathrm{omb}(X)}\\
&\text{(i)} && Z\ll X'\\
&\text{(i bis)} && \mathrm{omb}(Z)\subset\mathrm{omb}(X')\\
&\text{(ii)} && \mathrm{omb}(Z)^\circ\subset\mathrm{omb}(X') && \Downarrow\,\text{par ML 2.}\\
&\text{(iii)} && \mathrm{omb}(Z)^\circ\cap\mathrm{omb}(X')\neq\varnothing
\end{aligned}
\]

Il suffit de \unsure{noter|voir} $(\mathrm{iii})\Longrightarrow(\mathrm{i})$, et c'est ce qu'assure le lemme précédent.

\uline{Corollaire 3} \deleted{\unsure{Soit|Soient}}Dans les notations précédentes, $Z\ll X$, conditions équivalentes
\[
\begin{aligned}
&\text{(i)} && Z\overset{\circ}{\ll}X'\\
&\text{(ii)} && \mathrm{omb}(Z)^\circ\subset\mathrm{omb}(X')^\circ\\
&\text{(iii)} && \mathrm{omb}(Z)^\circ\cap\mathrm{omb}(X')^\circ\neq\varnothing\\
&\text{(iv)} && \deleted{\unsure{Z}}X'\ \text{est minimal dans }\widetilde X^{Z}=\{X'\in\widetilde X\mid Z\ll X'\}\\
&\deleted{\mathrm{(i)}\Rightarrow\mathrm{(ii)}\ \text{\unsure{car}}}\\
&\text{(\deleted{\unsure{i}}v)} && X'\ \text{est un plus petit \unsure{élémt|élément} dans }\widetilde X^{Z}
\end{aligned}
\]

On a
\diagram{\begin{tikzcd}[ampersand replacement=\&,row sep=small,column sep=small]
\& {\mathrm{(v)}}\arrow[d,Rightarrow] \& \&\\
{\mathrm{(i)}}\arrow[r,Leftrightarrow] \& {\mathrm{(iv)}} \& {\mathrm{(ii)}}\arrow[r,Rightarrow] \& {\mathrm{(iii)}}
\end{tikzcd}}
\unsure{trivialemt|trivialement}. On va prouver $(\mathrm{iii})\Longrightarrow(\mathrm{v})$ et $(\mathrm{i})\Longrightarrow(\mathrm{ii})$, ce qui suffira.

$(\mathrm{i})\Longrightarrow(\mathrm{ii})$ a été vu : si $x\in\mathrm{om}(Z)^\circ$ (\unsure{existe} par ML 2) \unsure{alors|i.e.} $x\overset{\circ}{\ll}Z$, et $Z\overset{\circ}{\ll}X'$, \unsure{on a} \unsure{vu|on} (Lemme 2) que $x\overset{\circ}{\ll}\deleted{\unsure{Z}}X'$ i.e $x\in\mathrm{omb}(X')^\circ$, donc (ii)

$(\mathrm{iii})\Longrightarrow(\mathrm{v})$ Soit $x\in\mathrm{omb}(Z)^\circ\cap\mathrm{omb}(X')^\circ$. Prouvons que $X'$ est plus petit dans $\widetilde X^{Z}$, i.e.
\scanpage{30}\mspage{29}

$Z\ll X''\leq X$, $\Longrightarrow X'\leq X''$. Soit donc \deleted{$x\in\mathrm{omb}(Z)^\circ$ (ML 2)} comme $x\in\mathrm{omb}(\deleted{\unsure{Z}}X')^\circ$, \deleted{\unsure{on voit que}} \inserted{$x\in\mathrm{omb}(X'')$, on} aura $\mathrm{omb}(X')^\circ\cap\mathrm{omb}(X'')\neq\varnothing$, d'où $X'\leq X''$ par lemme 1

OK.

\uline{Corollaire \deleted{\unsure{3}}\unsure{4}} \quad $X\overset{\circ}{\ll}Y\overset{\circ}{\ll}Z\Longrightarrow X\overset{\circ}{\ll}Z$

\margin{$\mathrm{At}_{M}\,5$}
\unsure{Immediat|Immédiat} sur le critère (ii) du cor. 3

Il est temps maintenant de faire intervenir la relation $\parallel$ sur \deleted{\unsure{$M$}}$L$. On va \unsure{maintenant} (je pense) utiliser le caractère minimal des $x\in L$ (?)

\deleted{\unsure{ML 1}} Je rappelle

(ML 0) la relation $x\parallel y$ \inserted{sur $L$} est symétrique et antiréflexive (elle implique $x\unsure{\neq|=}y$)

On pose dans $M$
\[
\begin{aligned}
X \mathrel{\tikz[baseline=-.39ex,x=1em,y=1em]{\draw[line width=.05em] (0,.55)--(.70,.27)--(0,0) (.70,.55)--(0,.27)--(.70,0);}} Y\quad &\overset{\mathrm{def}}{\Longleftrightarrow}\quad
\forall\begin{aligned}&x\in\mathrm{omb}(X)\smallsetminus\mathrm{omb}(\unsure{L|\llcorner})\\[-2pt]
 &y\in\mathrm{omb}(Y)\smallsetminus\mathrm{omb}(L)\end{aligned}
\quad\text{on a }x\parallel y\\
&\text{où }\mathrm{omb}(L)\overset{\mathrm{\unsure{defn|def}}}{=}
 \bigcup_{\substack{Z\in\widetilde X\cap\widetilde Y\\\text{i.e. }Z\leq X,Y}}\mathrm{omb}(Z)
\end{aligned}
\tag*{(1)}
\]
% slanted
\margin{relation symétrique, réflexive \deleted{elle est \unsure{antiréflexive} \unsure{puisque|qu'on}}.}

\diagram{\begin{tikzcd}[ampersand replacement=\&,row sep=small,column sep=small]
\mathrm{omb}(\unsure{X|x})\arrow[dr,dash] \& {\deleted{\mathrm{omb}(Y)}} \& \mathrm{omb}(Y)\arrow[dl,dash]\\
\& {\mathrm{omb}(X)\cap\mathrm{omb}(Y)}\arrow[d,dash] \&\\
\& \mathrm{omb}(L) \&
\end{tikzcd}}

\unsure{Donc} cela implique, par ML 0, que \deleted{$\mathrm{omb}(X)\cap\mathrm{omb}(Y)$}
\[
\mathrm{omb}(X)\cap\mathrm{omb}(Y)=\mathrm{omb}(L)
\]

% slanted
\margin{inutile ici, sera contenu dans \unsure{lemme} 5}
\uline{Lemme 3} Si $X,Y\leq Z$, on a $X \mathrel{\tikz[baseline=-.39ex,x=1em,y=1em]{\draw[line width=.05em] (0,.55)--(.70,.27)--(0,0) (.70,.55)--(0,.27)--(.70,0);}} Y$

Soit en effet \deleted{$x\in\mathrm{omb}$}
\scanpage{31}\mspage{30}

\diagram{\begin{tikzcd}[ampersand replacement=\&,row sep=0pt,column sep=1pt]
\& \& Z \& \&\\
\& \rotatebox{44}{$\leq$} \& \& \rotatebox{135}{$\leq$} \&\\
X \& \& \& \& Y\\
\& \rotatebox{135}{$\leq$} \& \& \rotatebox{44}{$\leq$} \&\\
\& \& {\unsure{L|\llcorner}} \& \&
\end{tikzcd}}

En effet, soit $x\in\mathrm{omb}(X)\smallsetminus\mathrm{omb}(L)$, je dis qu'on a $x\notin\mathrm{omb}(Y)$ -- en effet $\mathrm{omb}(X)\cap\mathrm{omb}(Y)=\mathrm{omb}(L)$ par ML 1 appliqué à $X,Y$. On utilisera \inserted{alors}

\uline{ML 4} Soit $Y\leq\deleted{Z}\unsure{X}$, \inserted{\deleted{\unsure{(ald}}(\unsure{ald} $Y\leq Z$)} et $x\in\mathrm{omb}(\deleted{\unsure{Z}}X)\smallsetminus\mathrm{omb}(\deleted{\unsure{X}}Y)$, (\unsure{donc|avec} $Y\neq\deleted{\unsure{Z}}X$ \ldots), $y\in\mathrm{omb}(Y)$, alors $x\parallel y$
\margin{$\mathrm{At}\,L\,3$}

\uline{Lemme} 4 Soient $X,Y\in M$, $X \mathrel{\tikz[baseline=-.39ex,x=1em,y=1em]{\draw[line width=.05em] (0,.55)--(.70,.27)--(0,0) (.70,.55)--(0,.27)--(.70,0);}} Y$. Alors l'énoncé du lemme 1 est valable pour $X,Y$.

% slanted
\margin{(N'utilise pas ML 4) C'est $\mathrm{At}_{M}\,6$}
Il suffit encore de prouver
\[
\mathrm{omb}(X)^\circ\cap\mathrm{omb}(Y)\neq\varnothing\quad\Longrightarrow\quad X\leq Y,
\]

\deleted{\unsure{Pour|Pren}} Soit donc
\[
x\in\mathrm{omb}(X)^\circ\cap\mathrm{omb}(Y)
\]
i.e.
\[
x\overset{\circ}{\ll}X,\quad x\ll Y
\]

\deleted{On ne peut avoir à la fois \unsure{$x\in\mathrm{omb}(X)$}} \deleted{$x\notin$}

En fait, comme $\mathrm{omb}(X)\cap\mathrm{omb}(Y)=\mathrm{omb}(L)$ \deleted{\unsure{on}} $\exists Z$ avec
\[
x\ll Z\leq X,Y
\]
comme $x\overset{\circ}{\ll}X$, on a $Z=X$, donc $X\leq Y$, qed.

\uline{Lemme} 5 Soit $X,Y,X',Y'$ avec \deleted{\unsure{$X \mathrel{\tikz[baseline=-.39ex,x=1em,y=1em]{\draw[line width=.05em] (0,.55)--(.70,.27)--(0,0) (.70,.55)--(0,.27)--(.70,0);}} Y$},} $X'\leq X$, $Y'\leq Y$, Alors $X \mathrel{\tikz[baseline=-.39ex,x=1em,y=1em]{\draw[line width=.05em] (0,.55)--(.70,.27)--(0,0) (.70,.55)--(0,.27)--(.70,0);}} Y\Longrightarrow X' \mathrel{\tikz[baseline=-.39ex,x=1em,y=1em]{\draw[line width=.05em] (0,.55)--(.70,.27)--(0,0) (.70,.55)--(0,.27)--(.70,0);}} Y'$.

Soit en effet
\[
\left\{\begin{aligned}
&x\in\mathrm{omb}(X')\smallsetminus\mathrm{omb}(L')\\
&y\in\mathrm{omb}(Y')\smallsetminus\mathrm{omb}(L')\\
&\qquad\text{avec }\mathrm{omb}\,L'\overset{\mathrm{\unsure{defn|def}}}{=}
   \bigcup_{Z\in\widetilde X'\cap\widetilde Y'}\mathrm{omb}(Z)
\end{aligned}\right.
\]
prouvons $x\parallel y$.
\scanpage{32}\mspage{31}
Si $x,y\notin\mathrm{omb}(\unsure{L|\llcorner})$, cela résulte de l'\unsure{hypothèse} $X\mathrel{\tikz[baseline=-.39ex,x=1em,y=1em]{\draw[line width=.05em] (0,.55)--(.70,.27)--(0,0) (.70,.55)--(0,.27)--(.70,0);}}Y$. \unsure{Sinon|Si non} \deleted{\unsure{$x\notin$|$x\in$}} \unsure{)|,} soit p.ex. $x\in\mathrm{omb}(L)$, \deleted{Pour} i.e $\exists\,Z\in\widetilde X\cap\widetilde Y$ (i.e. $\unsure{2|Z}\leq X,Y$) avec $x\ll Z$

\unsure{[}Appliquons $\mathrm{ML}\ 1$ à
\diagram{\begin{tikzcd}[ampersand replacement=\&,column sep=.15em,row sep=.12em]
X' \& \leq \& X \\
\rotatebox{90}{$\ll$} \& \& \rotatebox{90}{$\leq$} \\
x \& \ll \& Z
\end{tikzcd}}
, on trouve \deleted{$S\in M$ \unsure{--} $S\in\widetilde X$ avec.}
\diagram{\begin{tikzcd}[ampersand replacement=\&,column sep=.15em,row sep=.13em]
\& \& X' \& \leq \& X \\
\& \& \rotatebox{90}{$\leq$} \& \& \rotatebox{90}{$\leq$} \\
\& \& S \& \leq \& Z \& (\leq Y)\\
\& \rotatebox{45}{$\ll$} \& \& \& \& \\
x \& \& \& \& \&
\end{tikzcd}}
i.e.
\diagram{\begin{tikzcd}[ampersand replacement=\&,column sep=.1em,row sep=.12em]
X' \& \& \& \& Y\\
\& \rotatebox{135}{$\leq$} \& \& \rotatebox{45}{$\leq$} \& \\
\& \& S \& \& \\
\& \rotatebox{45}{$\ll$} \& \& \& \\
x \& \& \& \&
\end{tikzcd}}
\deleted{donc $S\in\widetilde L\overset{\mathrm{def}}{=}\widetilde X\cap\widetilde Y$]}

\deleted{donc $x\ll$ \unsure{Donc}} \inserted{On a} $x\ll Y$ \inserted{i.e $x\in\mathrm{omb}(Y)$}, et pour prouver \deleted{Je \unsure{dis|di}} $x\parallel y$, il suffit \uline{par} \uuline{$\mathrm{ML}\ 4$} de prouver $x\notin\mathrm{omb}(Y')$. \deleted{i.e. $x\unsure{<|\ll}$} \inserted{Mais si} on avait $x\in\mathrm{omb}(Y')$, on aurait par $\mathrm{ML}\ 1$ appliqué à
\diagram{\begin{tikzcd}[ampersand replacement=\&,column sep=.15em,row sep=.14em]
S \& \leq \& Y\\
\deleted{\rotatebox{90}{$\unsure{\leq|<}$}}\rotatebox{90}{$\ll$} \& \& \rotatebox{90}{$\leq$}\\
x \& \ll \& Y'
\end{tikzcd}}
un $T$ \deleted{\unsure{$\in$}} tel que
\[
 x\ll T\leq S,Y'
\]
et on aurait donc aussi
\[
 T\leq X',Y'\qquad\text{puisque } S\leq X'
\]
donc $x\in\mathrm{omb}(L')$, contrairement à l'hyp, qed.

\deleted{\uline{Lemme \unsure{4|6}}} Il faut prouver maintenant seulement $\mathrm{At}_M\,4$ \deleted{$\mathrm{At}_M\,\unsure{5|3}$} b) ($\mathrm{At}_M\,5,6$ et $\mathrm{At}_M\,4$ a) sont déjà prouvés) pour conclure que $(M,\leq,\ll,\underset{M}{\mathrel{\tikz[baseline=-.39ex,x=1em,y=1em]{\draw[line width=.05em] (0,.55)--(.70,.27)--(0,0) (.70,.55)--(0,.27)--(.70,0);}}})$ (et \deleted{la \unsure{famille|famil}} $\unsure{\mathcal F|\mathscr F|F}={}$ens des parties fermées de $M$ \unsure{formés|formées} d'él. $2$ à $2$ compatibles pour $\underset{M}{\mathrel{\tikz[baseline=-.39ex,x=1em,y=1em]{\draw[line width=.05em] (0,.55)--(.70,.27)--(0,0) (.70,.55)--(0,.27)--(.70,0);}}}$)
\scanpage{33}\mspage{32}
\deleted{\unsure{forment un}} correspondent \unsure{:|à} un atelier $\mathrm{At}\ 1$ -- $\mathrm{At}\ 7$. Donc c'est le contenu du

\uline{Lemme} 6 \quad Soient $X,Y,\unsure{2|Z}$ avec
\[
 Z\overset{\circ}{\ll}X,\ Z\overset{\circ}{\ll}Y,\quad X\mathrel{\tikz[baseline=-.39ex,x=1em,y=1em]{\draw[line width=.05em] (0,.55)--(.70,.27)--(0,0) (.70,.55)--(0,.27)--(.70,0);}}Y
\]
alors $X=Y$.

Ça résultera du cor. plus général

\[
\left|\begin{aligned}
&\text{\uline{Cor} \quad On a, si }X\mathrel{\tikz[baseline=-.39ex,x=1em,y=1em]{\draw[line width=.05em] (0,.55)--(.70,.27)--(0,0) (.70,.55)--(0,.27)--(.70,0);}}Y\\
& Z\overset{\circ}{\ll}X,\quad Z\ll Y,\Longrightarrow X\leq\deleted{\unsure{<|\ll}}Y
\end{aligned}\right.
\]
\uline{Dém} \quad Soit $x\in\mathrm{omb}(Z)^\circ$ (ML 2). On a donc \deleted{$x\ll$ \unsure{om|ou|on}}
\[
 x\overset{\circ}{\ll}Z\overset{\circ}{\ll}X
\]
donc (lemme \unsure{2|?}, cor. 4 p.29) $x\overset{\circ}{\ll}X$, donc $x\in\mathrm{omb}(X)^\circ\cap\mathrm{omb}(Y)$, lequel est $\neq\varnothing$, donc $X\leq Y$ (lemme 4).

Il faut voir que la relation de \unsure{disjoint|disjont} $\parallel$ sur $L$ est bien la relation de disjonction au sens de l'atelier qu'on vient de construire, i.e.

\uline{Lemme} 7 \quad \deleted{\unsure{Soit|So}} a) Soient $x,y\in L$. Alors $x\mathrel{\tikz[baseline=-.39ex,x=1em,y=1em]{\draw[line width=.05em] (0,.55)--(.70,.27)--(0,0) (.70,.55)--(0,.27)--(.70,0);}}y$ \unsure{sss|ssi} $x=y$ ou $x\parallel y$.

b) On a $x\parallel y$ sss $x\mathrel{\tikz[baseline=-.39ex,x=1em,y=1em]{\draw[line width=.05em] (0,.55)--(.70,.27)--(0,0) (.70,.55)--(0,.27)--(.70,0);}}y$ et $\mathrm{Inf}^{\deleted{\unsure{L|F}}\unsure{\mathcal F|\mathscr F|F},\leq}(x,y)=\varnothing_{\unsure{F|\mathcal F}}$. (i.e. $x$ et $y$ disjoints au sens de l'atelier)

\unsure{Pour|Dans} a, il suffit de voir le cas $x\neq y$. À cause du \uline{caractère minimal} de $x,y$, on aura \inserted{$\mathrm{omb}\,x=\{x\},\ \mathrm{omb}(y)=\{y\},$} $\unsure{L|\llcorner}=\varnothing$, et le critère \deleted{est bien} par défaut $x\mathrel{\tikz[baseline=-.39ex,x=1em,y=1em]{\draw[line width=.05em] (0,.55)--(.70,.27)--(0,0) (.70,.55)--(0,.27)--(.70,0);}}y$ est bien $x\parallel y$. L'énoncé b) \unsure{suit;|suit,}
